% 2026-07-01 OmegaSim Feedback Revision
% Addresses: cognitive appraisal, latent motivation, semantic/artifact fields,
% costly prediction, hysteresis, residual-state lobe discovery, linear controls,
% and experiment families A6--A8.

\section*{Revision Addendum (2026-07-01): Cognitive Appraisal, Residual Lobes, and New Experiment Families}

This addendum responds to the 2026-07-01 feedback on the original note.  It
preserves all prior material and adds eight substantive extensions.

\subsection*{R1. Logistic Response as Bounded Thresholded Cognitive Appraisal}

The original note used a generic saturating nonlinearity $S$.  We now interpret
the nonlinear response as a \emph{cognitive appraisal function}: each agent or
hive component evaluates incoming signals through a bounded, thresholded
transform---not as arbitrary chaos injection, but as a principled encoding of
attention allocation, urgency, and resource commitment.

\begin{definition}[Cognitive appraisal response]
The componentwise response is a logistic (sigmoid) function
\[
  \sigma_\alpha(x;\theta,\beta)=\frac{1}{1+\exp(-\alpha(x-\theta))},
\]
where $\alpha>0$ is the appraisal steepness (sensitivity to deviations from
threshold $\theta$) and $\beta$ is the saturation level (asymptotic maximum
response).  The effective response in the hive update replaces $S(\cdot)$ with
\[
  S_i(\cdot)\;\mapsto\; \beta_i\,\sigma_{\alpha_i}(w_i^\top x_k + b_i;\theta_i,\beta_i),
\]
giving each component $i$ its own threshold, steepness, and amplitude.
\end{definition}

\begin{proposition}[Appraisal threshold as bifurcation coordinate]
The appraisal steepness $\alpha$ acts as a gain parameter: low $\alpha$ yields
near-linear, contractive dynamics; high $\alpha$ creates a switching regime with
sharp transitions.  The threshold $\theta$ shifts the operating point.  Together
$(\alpha,\theta)$ provide a two-parameter family of bifurcation diagrams that
interpolates between smooth contraction and thresholded multi-state dynamics.
\end{proposition}

\begin{proof}
For $\alpha\to 0$, $\sigma_\alpha$ approaches a constant linear function with
slope $\beta/4$, recovering the low-gain contractive regime.  For $\alpha\to\infty$,
$\sigma_\alpha$ approaches a Heaviside step, introducing discontinuity-class
bifurcations.  Intermediate $\alpha$ produces smooth but steep transitions whose
Jacobian norm scales as $\alpha\beta/4$ at the threshold, directly controlling
whether the lifted map is contractive.
\end{proof}

\subsection*{R2. Latent Motivational State with Explicit Dynamics}

\begin{definition}[Latent motivational state]
Introduce a latent motivational vector $\mu_k\in\mathbb R^p$ with explicit
dynamics
\[
  \mu_{k+1}=(1-\lambda)\mu_k+\lambda\,\Phi(x_k,\mu_k)+\zeta_k,
\]
where $\lambda\in(0,1)$ is the motivation update rate, $\Phi$ is a feedback
function from current hive state and current motivation to updated motivation,
and $\zeta_k$ is endogenous drive noise.  The motivation $\mu_k$ modulates the
appraisal thresholds:
\[
  \theta_i(\mu_k)=\theta_i^{(0)}+\gamma_i^\top\mu_k,
\]
so that motivational state shifts what deviations are attended to.
\end{definition}

\begin{assumption}[Motivational timescale separation]
We assume $\lambda\ll 1/T_{\rm decision}$, i.e., motivational state evolves
slower than the decision update.  This yields a slow--fast system where $x_k$
is the fast variable and $\mu_k$ is the slow variable, enabling
singular-perturbation analysis.
\end{assumption}

\begin{proposition}[Slow motivation induces bifurcation sweeping]
As $\mu_k$ evolves slowly, the appraisal thresholds $\theta_i(\mu_k)$ drift,
causing the fast dynamics' bifurcation parameters to sweep.  This can produce
intermittent visits to multiple attractor lobes, hysteresis across bifurcation
boundaries, and crisis-induced switching---all driven by internal motivational
state rather than external forcing.
\end{proposition}

\subsection*{R3. Separate Semantic and Artifact Fields and Their Coupling}

\begin{definition}[Semantic field]
The semantic field $s_k\in\mathbb R^{d_s}$ captures the representational content
flowing through the hive: message semantics, task descriptions, shared memory
content, and interpretive type assignments.  Its dynamics are
\[
  s_{k+1}=\Psi_s(s_k,x_k,\mu_k)+\xi^s_k,
\]
where $\Psi_s$ may involve embedding-space operations, semantic similarity
weighting, and type-dependent routing.
\end{definition}

\begin{definition}[Artifact field]
The artifact field $a_k\in\mathbb R^{d_a}$ tracks produced artifacts: code,
documents, test results, structured outputs.  Its dynamics are
\[
  a_{k+1}=\Psi_a(a_k,s_k,x_k,\mu_k)+\xi^a_k,
\]
where $\Psi_a$ includes quality evaluation, version tracking, and resource cost
accounting.
\end{definition}

\begin{definition}[Field coupling]
The coupled system replaces the monolithic state $x_k$ with
\[
  \tilde x_k=(x_k^{\rm op},s_k,a_k,\mu_k),
\]
where $x_k^{\rm op}$ are operational/coordination variables.  The semantic field
influences operational dynamics through appraisal thresholds $\theta_i(s_k)$,
and the artifact field feeds back through resource load $r_k=g(a_k)$ and
quality metrics $q_k=Q(a_k)$.  Bidirectional coupling strength is parameterized
by $\chi_{sa}$ (semantic$\to$artifact) and $\chi_{as}$ (artifact$\to$semantic).
\end{definition}

\begin{designprinciple}[Separate fields enable distinct diagnostics]
Treating $s_k$ and $a_k$ as distinct dynamical fields allows separate Lyapunov
analysis, attractor characterization, and surrogate controls for each.  Semantic
complexity without artifact complexity, or vice versa, is diagnostically distinct
from joint complexity.
\end{designprinciple}

\subsection*{R4. Costly Prediction and Resource Accounting}

\begin{definition}[Prediction cost]
Each predictive step---whether by an agent predicting outcomes or by the hive
forecasting its own state---incurs a resource cost
\[
  c_k^{\rm pred}=c_{\rm fixed}+c_{\rm var}\|p_k\|^2,
\]
where $p_k$ is the prediction depth or model complexity and $c_{\rm var}>0$.
Total resource consumption must satisfy
\[
  \sum_{i} c_{k,i}^{\rm pred}+c_{k,i}^{\rm act}\le R_{\max},
\]
where $R_{\max}$ is the hive-wide resource budget per timestep.
\end{definition}

\begin{definition}[Resource-constrained dynamics]
The hive update with resource accounting becomes
\[
  x_{k+1}=S\bigl(W_0x_k+W_1x_{k-d}+W_mM_k+B u_k+\xi_k;\,R_{\max}-c_k^{\rm pred}\bigr),
\]
where the remaining budget $R_{\max}-c_k^{\rm pred}$ gates the effective
amplitude of activation.  When prediction costs are high, effective gain drops,
pushing the system toward the contractive regime; when prediction is cheap,
full nonlinear dynamics operate.
\end{definition}

\begin{proposition}[Cost-induced self-regulation]
If prediction cost scales with the complexity of the dynamics (e.g.,
$c_k^{\rm pred}\propto$ local finite-time Lyapunov exponent), then highly
chaotic regimes automatically incur higher costs, reducing effective gain and
creating a self-regulating feedback loop.  This provides an intrinsic mechanism
for bounded complexity without external clipping.
\end{proposition}

\subsection*{R5. Hysteresis and Adaptive Thresholds}

\begin{definition}[Adaptive appraisal thresholds]
Appraisal thresholds evolve according to
\[
  \theta_{i,k+1}=\theta_{i,k}+\eta_i\bigl(\sigma_{\alpha_i}(w_i^\top x_k;\theta_{i,k})-p_i^{\ast}\bigr),
\]
where $p_i^{\ast}\in(0,1)$ is a target activation rate and $\eta_i>0$ is the
adaptation rate.  This drives each component toward a target operating point,
implementing homeostatic regulation.
\end{definition}

\begin{proposition}[Adaptive thresholds create hysteresis]
When the adaptation rate $\eta_i$ is small but nonzero, the threshold
$\theta_{i,k}$ tracks the recent activation level.  On upward parameter sweeps
(increasing gain $\alpha$), the threshold rises to suppress over-activation; on
downward sweeps, it lowers slowly.  This asymmetric adaptation produces
hysteresis: the system's state at gain $\alpha$ on the upward sweep differs from
its state at the same $\alpha$ on the downward sweep.  The hysteresis loop width
scales as $O(\eta_i / \lambda_{\rm bif})$ where $\lambda_{\rm bif}$ is the
bifurcation parameter sweep rate.
\end{proposition}

\begin{proof}
Let $\bar\theta(\alpha)$ be the quasi-steady threshold satisfying
$\sigma_{\alpha}(w^\top x^\ast(\alpha,\bar\theta);\bar\theta)=p^\ast$.
Under slow sweep $\alpha_k=\alpha_0+\epsilon k$, the actual threshold lags:
$\theta_{k}\approx\bar\theta(\alpha_k)-\frac{\eta_i}{\epsilon}\bar\theta'(\alpha_k)$.
The correction term is direction-dependent (it changes sign with sweep
direction), producing a hysteresis gap proportional to $\eta_i/\epsilon$.
\end{proof}

\subsection*{R6. Residual-State Lobe Discovery}

\begin{definition}[Residual state]
Define the residual state as the deviation of the full system state from the
best-fit low-dimensional attractor model:
\[
  \rho_k = \tilde x_k - \Pi_{\mathcal A}(\tilde x_k),
\]
where $\tilde x_k=(x_k^{\rm op},s_k,a_k,\mu_k)$ is the full coupled state,
$\Pi_{\mathcal A}$ is the projection onto the fitted attractor model
$\mathcal A$ (e.g., a principal manifold, delay-coordinate embedding, or
low-order DMD approximation), and $\rho_k$ is the residual.  The residual
captures structure not explained by the low-dimensional attractor model.
\end{definition}

\begin{definition}[Residual lobe]
A residual lobe is a persistent cluster in residual-state space
$\{\rho_k\}$ identified by a clustering algorithm (e.g., Gaussian mixture,
HDBSCAN, or $k$-medoids in delay-coordinate space) with the following
properties:
\begin{enumerate}
  \item \emph{Persistence}: the cluster survives across multiple validation
        windows of length $\ge 10\,T_{\rm relax}$.
  \item \emph{Separation}: the cluster's centroid is at least $2\sigma$ from
        other cluster centroids in residual norm.
  \item \emph{Reproducibility}: the cluster appears in $\ge 80\%$ of
        independent runs with different initial histories under the same
        parameter setting.
  \item \emph{Semantic distinguishability}: Hyperseed event records for
        trajectories within the lobe differ systematically (in agent roles,
        artifact types, or semantic content) from records for other lobes.
\end{enumerate}
\end{definition}

\begin{designprinciple}[Lobe discovery as a validation tool]
Residual lobe discovery provides a concrete method for detecting whether the
full system has multi-lobe attractor structure that is invisible to the
low-dimensional observable.  If lobes are semantically distinguishable, they
identify distinct cognitive regimes (e.g., exploration vs.\ exploitation,
problem-attacking vs.\ reflective) that the hive transitions between.
\end{designprinciple}

\subsection*{R7. Amplitude-Matched Linear Controls}

\begin{definition}[Amplitude-matched linear control]
Given a nonlinear hive system $x_{k+1}=G(x_k,\ldots,x_{k-m})$ with operational
state norm $\|x_k\|\le X_{\max}$ in the regime of interest, the
amplitude-matched linear control is
\[
  x_{k+1}^{\rm lin}=A_0 x_k + A_1 x_{k-d} + B u_k + \bar c,
\]
where $A_0,A_1$ are obtained by linearizing $G$ around the regime mean
$\bar x = \frac{1}{T}\sum_{k} x_k$ and $\bar c$ is chosen so that
$\|x_k^{\rm lin}\|\approx\|x_k\|$ in the matched regime.  The linear control
uses the same noise path $\xi_k$ (common random numbers) and the same delay
structure as the nonlinear system.
\end{definition}

\begin{proposition}[Linear control as a surrogate baseline]
The amplitude-matched linear control isolates the contribution of nonlinearity
from delay and noise.  If the linear control exhibits comparable spectral
breadth, Lyapunov exponent, or recurrence statistics, the nonlinear system's
complexity is attributable to delay/noise, not to nonlinear structure.  Only
when the nonlinear system significantly exceeds the linear control on these
metrics can nonlinear strange-attractor claims be supported.
\end{proposition}

\begin{proof}
The linear control shares delay structure, noise realization, and amplitude
with the nonlinear system; it differs only in the replacement of $G$ by its
linearization.  Therefore any excess complexity in the nonlinear system is
attributable to the nonlinear terms.  This is a standard matched-control
argument from the design-of-experiments tradition.
\end{proof}

\subsection*{R8. Experiment Families A6, A7, and A8}

\subsubsection*{A6: Cognitive Appraisal Bifurcation Scan}

\textbf{Hypothesis:} Varying appraisal steepness $\alpha$ and threshold
$\theta$ across the hive produces a smooth bifurcation from contractive
fixed-point dynamics through damped oscillation to multi-lobe recurrent
dynamics, mediated by the cognitive appraisal function rather than by external
noise injection.

\textbf{Sweep variables:}
\begin{itemize}
  \item Appraisal steepness $\alpha\in\{0.5,1,2,5,10,20\}$
  \item Threshold $\theta\in\{-2,-1,0,1,2\}$ (relative to baseline input
        distribution)
  \item Motivational coupling $\gamma\in\{0,0.1,0.5,1.0\}$
\end{itemize}

\textbf{Controls:}
\begin{itemize}
  \item Amplitude-matched linear control (R7) at each $(\alpha,\theta)$ point
  \item No-motivational-dynamics control ($\gamma=0$, frozen $\mu_k$)
  \item Phase-randomized surrogate
\end{itemize}

\textbf{Acceptance criteria:}
\begin{enumerate}
  \item Bifurcation diagram in $(\alpha,\theta)$ shows a qualitative
        transition from fixed-point to oscillatory to multi-lobe dynamics.
  \item The transition region has positive finite-time Lyapunov exponent
        while the linear control does not.
  \item Motivational coupling ($\gamma>0$) widens the multi-lobe region
        compared to $\gamma=0$.
  \item Hysteresis is observable on upward vs.\ downward $\alpha$ sweeps with
        adaptive thresholds.
\end{enumerate}

\subsubsection*{A7: Residual Lobe Discovery and Validation}

\textbf{Hypothesis:} The full coupled system $(x^{\rm op},s,a,\mu)$ has
multi-lobe attractor structure that is partially invisible to the low-dimensional
observable $O(E_{\le t})$.  Residual lobes are semantically distinguishable.

\textbf{Sweep variables:}
\begin{itemize}
  \item Attractor model dimension $d_{\mathcal A}\in\{2,3,5,8\}$
  \item Clustering algorithm: HDBSCAN, GMM, $k$-medoids
  \item Delay embedding dimension for residual space: $m_{\rm emb}\in\{3,5,10\}$
  \item Validation window length: $\{5,10,20\}\cdot T_{\rm relax}$
  \item Number of independent runs: $\ge 20$ per parameter setting
\end{itemize}

\textbf{Controls:}
\begin{itemize}
  \item Random-label surrogate: cluster assignments randomly permuted across
        time to test whether semantic distinguishability exceeds chance.
  \item Single-field controls: residual computed using only $x^{\rm op}$, only
        $s$, only $a$, or only $\mu$, to test which field's residual carries
        lobe structure.
  \item Null-model residual: residual computed against a trivial (zero-order)
        attractor model, i.e., the mean state, to verify that the clustering
        captures genuine attractor structure and not just noise.
\end{itemize}

\textbf{Acceptance criteria:}
\begin{enumerate}
  \item At least one residual lobe satisfies all four properties from the
        residual lobe definition (R6): persistence, separation, reproducibility,
        and semantic distinguishability.
  \item The random-label surrogate produces $<5\%$ false-positive lobe rate
        at the same clustering thresholds.
  \item Single-field controls identify which field(s) contribute most to lobe
        structure, enabling targeted interpretation.
  \item Lobe membership predicts subsequent behavioral regime with above-chance
        accuracy ($p<0.01$ by permutation test).
\end{enumerate}

\subsubsection*{A8: Costly Prediction and Self-Regulation Experiment}

\textbf{Hypothesis:} When prediction cost scales with dynamical complexity
(high local Lyapunov exponent $\Rightarrow$ higher $c_k^{\rm pred}$), the hive
system self-regulates into a bounded-complexity regime: chaotic episodes are
automatically damped by resource depletion, and the system settles into an
intermittent regime alternating between high-complexity bursts and
resource-recovery quiescence.

\textbf{Sweep variables:}
\begin{itemize}
  \item Prediction cost coefficient $c_{\rm var}\in\{0,0.01,0.1,0.5,1.0\}$
  \item Resource budget $R_{\max}\in\{0.5,1,2,5\}$ (relative to median
        unconstrained activation cost)
  \item Cost complexity coupling: $c_k^{\rm pred}\propto |\lambda_{\rm FTLE}|^p$
        with exponent $p\in\{0,1,2\}$ (no coupling, linear, quadratic)
  \item Prediction depth $p_k\in\{1,2,5,10\}$ steps ahead
\end{itemize}

\textbf{Controls:}
\begin{itemize}
  \item Zero-cost control ($c_{\rm var}=0$): standard dynamics without
        resource accounting, to establish the unconstrained complexity baseline.
  \item Fixed-cost control: $c_k^{\rm pred}=c_{\rm fixed}$ (constant regardless
        of complexity), to distinguish resource scarcity effects from
        complexity-coupled self-regulation.
  \item Amplitude-matched linear control (R7) with the same resource budget,
        to separate nonlinear self-regulation from linear resource gating.
\end{itemize}

\textbf{Acceptance criteria:}
\begin{enumerate}
  \item With cost-complexity coupling ($p>0$), the system exhibits
        intermittent complexity: episodes of positive finite-time Lyapunov
        exponent alternated with quiescent recovery periods.
  \item The zero-cost control shows sustained chaos (no self-damping),
        confirming that the intermittent pattern is caused by cost coupling,
        not by other parameters.
  \item The fixed-cost control shows uniform suppression (no intermittency),
        confirming that the intermittency requires complexity-dependent cost,
        not just resource scarcity.
  \item Long-run complexity (time-averaged Lyapunov exponent) decreases
        monotonically with $c_{\rm var}$ for fixed $R_{\max}$, demonstrating
        a tunable self-regulation knob.
  \item Hysteresis in the complexity--cost relationship is observable when
        $c_{\rm var}$ is swept up and down, due to adaptive thresholds (R5).
\end{enumerate}

\subsection*{Summary of revisions}

This addendum adds: (R1) cognitive appraisal interpretation of the logistic
response, (R2) latent motivational state with slow--fast dynamics and
threshold modulation, (R3) separate semantic and artifact fields with
bidirectional coupling, (R4) costly prediction with resource accounting and
self-regulating complexity, (R5) adaptive thresholds producing hysteresis,
(R6) residual-state lobe discovery with clustering and validation, (R7)
amplitude-matched linear controls isolating nonlinear contributions, and (R8)
three new experiment families (A6, A7, A8) with explicit hypotheses, sweep
variables, controls, and acceptance criteria.  All original material is
preserved; this content is appended via \texttt{\\input} before
\texttt{\\end{document}}.